\documentclass[preview]{standalone}

\usepackage{amsmath}
\usepackage{amssymb}
\usepackage{stellar}
\usepackage{definitions}
\usepackage{bettelini}

\begin{document}

\id{probability-conditional-probability}
\genpage

\section{Conditional probability}

\begin{snippetdefinition}{conditional-probability-definition}{Conditional probability}
    Let \((\Omega,\mathcal F,\mathbb P)\) be a \probabilityspace and let
    \(A,B\in\mathcal F\). If \(\mathbb P(B)>0\), the
    \emph{conditional probability} of \(A\) given \(B\) is
    \[
        \mathbb P(A\mid B)
        \triangleq
        \frac{\mathbb P(A\intersection B)}{\mathbb P(B)}.
    \]
    If \(\mathbb P(A\mid B)>\mathbb P(A)\), then \(B\) positively conditions
    \(A\). If \(\mathbb P(A\mid B)<\mathbb P(A)\), then \(B\) negatively
    conditions \(A\). If \(\mathbb P(A\mid B)=\mathbb P(A)\), then \(B\)
    does not change the probability of \(A\).
\end{snippetdefinition}

\begin{snippetnote}{conditioning-on-null-events}{Conditioning on null events}
    Let \((\Omega,\mathcal F,\mathbb P)\) be a \probabilityspace and let
    \(A,B\in\mathcal F\). Ordinary \conditionalprobability given the \event
    \(B\) is defined only when \(\mathbb P(B)>0\).
    If one tries to define \(\mathbb P(A\mid B)\) as a number \(x\in[0,1]\)
    satisfying
    \[
        \mathbb P(A\intersection B)=x\mathbb P(B),
    \]
    then for \(\mathbb P(B)=0\) the equation has no unique solution.
\end{snippetnote}

\begin{snippetproposition}{conditional-probability-measure}{Conditional probability is a probability measure}
    Let \((\Omega,\mathcal F,\mathbb P)\) be a \probabilityspace and let
    \(B\in\mathcal F\) with \(\mathbb P(B)>0\).
    Define
    \[
        \mathbb P_B\colon\mathcal F\fromto[0,1],
        \qquad
        \mathbb P_B(A)=\mathbb P(A\mid B).
    \]
    Then \(\mathbb P_B\) is a \probabilitymeasure on \((\Omega,\mathcal F)\).
\end{snippetproposition}

\begin{snippetproof}{conditional-probability-measure-proof}{conditional-probability-measure}{Conditional probability is a probability measure}
    First,
    \[
        \mathbb P_B(\Omega)
        =
        \frac{\mathbb P(\Omega\intersection B)}{\mathbb P(B)}
        =
        \frac{\mathbb P(B)}{\mathbb P(B)}
        =
        1.
    \]
    Let \(\{A_n\}_{n\in\naturalnumbers}\subseteq\mathcal F\) be pairwise
    disjoint. Then the \event[events] \(A_n\intersection B\) are pairwise
    disjoint, and therefore
    \begin{align*}
        \mathbb P_B\left(\bigcup_{n=1}^\infty A_n\right)
        &=
        \frac{
            \mathbb P\left(
                \left(\bigcup_{n=1}^\infty A_n\right)\intersection B
            \right)
        }{\mathbb P(B)}
        =
        \frac{
            \mathbb P\left(
                \bigcup_{n=1}^\infty (A_n\intersection B)
            \right)
        }{\mathbb P(B)}
        \\
        &=
        \frac{
            \sum_{n=1}^\infty\mathbb P(A_n\intersection B)
        }{\mathbb P(B)}
        =
        \sum_{n=1}^\infty\mathbb P(A_n\mid B)
        =
        \sum_{n=1}^\infty\mathbb P_B(A_n).
    \end{align*}
\end{snippetproof}

\begin{snippetdefinition}{sample-space-partition-definition}{Sample space partition}
    Let \((\Omega,\mathcal F)\) be a \measurablespace.
    A finite or countable family
    \[
        \{B_i\}_{i\in I}\subseteq\mathcal F
    \]
    is a \emph{partition} of the \samplespace \(\Omega\) if
    \begin{enumerate}
        \item \(\Omega=\bigcup_{i\in I}B_i\);
        \item \(B_i\intersection B_j=\emptyset\) whenever \(i\neq j\).
    \end{enumerate}
\end{snippetdefinition}

\begin{snippettheorem}{total-probability-law}{Law of total probability}
    Let \((\Omega,\mathcal F,\mathbb P)\) be a \probabilityspace.
    Let \(\{B_i\}_{i\in I}\subseteq\mathcal F\) be a finite or countable
    \samplespacepartition of \(\Omega\), and assume that
    \(\mathbb P(B_i)>0\) for every \(i\in I\).
    Then, for every \(A\in\mathcal F\),
    \[
        \mathbb P(A)
        =
        \sum_{i\in I}\mathbb P(A\mid B_i)\mathbb P(B_i).
    \]
\end{snippettheorem}

\begin{snippetproof}{total-probability-law-proof}{total-probability-law}{Law of total probability}
    Since \(\{B_i\}_{i\in I}\) is a \samplespacepartition of \(\Omega\),
    the \event[events] \(A\intersection B_i\) are pairwise disjoint and
    \[
        A=A\intersection\Omega
        =
        A\intersection\left(\bigcup_{i\in I}B_i\right)
        =
        \bigcup_{i\in I}(A\intersection B_i).
    \]
    Hence, by countable additivity,
    \[
        \mathbb P(A)
        =
        \sum_{i\in I}\mathbb P(A\intersection B_i)
        =
        \sum_{i\in I}\mathbb P(A\mid B_i)\mathbb P(B_i).
    \]
\end{snippetproof}

\begin{snippettheorem}{bayes-theorem}{Bayes' theorem}
    Let \((\Omega,\mathcal F,\mathbb P)\) be a \probabilityspace.
    Let \(\{B_i\}_{i\in I}\subseteq\mathcal F\) be a finite or countable
    \samplespacepartition of \(\Omega\), and assume that
    \(\mathbb P(B_i)>0\) for every \(i\in I\).
    If \(A\in\mathcal F\) and \(\mathbb P(A)>0\), then for every \(k\in I\),
    \[
        \mathbb P(B_k\mid A)
        =
        \frac{\mathbb P(A\mid B_k)\mathbb P(B_k)}
        {\sum_{i\in I}\mathbb P(A\mid B_i)\mathbb P(B_i)}.
    \]
\end{snippettheorem}

\begin{snippetproof}{bayes-theorem-proof}{bayes-theorem}{Bayes' theorem}
    By the definition of \conditionalprobability,
    \[
        \mathbb P(B_k\mid A)
        =
        \frac{\mathbb P(A\intersection B_k)}{\mathbb P(A)}
        =
        \frac{\mathbb P(A\mid B_k)\mathbb P(B_k)}{\mathbb P(A)}.
    \]
    By the \totalprobabilitylaw,
    \[
        \mathbb P(A)
        =
        \sum_{i\in I}\mathbb P(A\mid B_i)\mathbb P(B_i),
    \]
    and substituting this expression for \(\mathbb P(A)\) gives the formula.
\end{snippetproof}

\end{document}
